\section{Deciding closed blocks of collapse arithmetic}\label{sec:arithmetic}

We now combine order comparison with the modular tables. A finite circuit
may use natural inputs, natural-valued affine gates, fixed iterates of
\(f\), entry norms, interval counts, and fixed-head collapses. Each gate
and call carries its true domain; in particular an affine gate's output
is required to be natural. Signed linear combinations remain unrestricted
in the final inequalities, equalities and congruences. The total shifted function
\(n\mapsto J_\lambda(N(\lambda)+n)\) is a convenient way to avoid
partial raw inputs. Alternatively, an explicitly specified computable
finite extension below \(N(\lambda)\) can be handled as finite cases.

\begin{theorem}[Closed-block decision]\label{thm:arithmetic}
For either control of Section~\ref{sec:controls}, closed existential
and closed universal Boolean combinations of linear inequalities,
equalities and congruences over finite circuits of fixed-head collapses
are decidable. The procedure preserves shared inputs and shared circuit
nodes. The fixed heads may be arbitrary ordinals below \(\varepsilon_0\).
\end{theorem}

The descent notation below uses nonzero limit heads. This loses no fixed
head in the theorem: write any infinite \(\alpha\) uniquely as
\(\lambda+k\), with \(\lambda\) a nonzero limit and \(k<\omega\),
and absorb \(k\) into the finite input of
\(\psi_f(\alpha+u)=\psi_f(\lambda+(k+u))\).
For finite \(\alpha\), the value is simply \(\alpha+u\).

The proof is constructive, but makes no useful running-time claim. Its
two termination measures are the number of raw coordinates and a finite
charge on the active count occurrences.

\subsection{The finite matrix language}

Flatten a finite circuit once, assigning one natural coordinate to each
node and conjoining its exact graph equation. A collapse gate is
\[
 y=Z_\lambda(t)+D_\lambda(t),\qquad t\ge N(\lambda).
\]
Affine gates give affine equalities. Every valuation in the circuit's
original domain has exactly one extension to these graph coordinates.
This remains true when the
original output test is negated. All graph equations stay conjoined
outside that test.

After flattening, every nonlinear function has one raw coordinate as its
argument. Fixed argument shifts are represented by affine aliases. The
ordinary nonlinear value fields are
\[
 E_P(x),\quad f^a(E_P(x))\ (a\ge1),\quad f^a(x)\ (a\ge1),
\]
and the count fields are \(\tau_P(x)\). Coefficients in rows and
congruences are fixed integers. An uniterated raw coordinate is linear,
not a separated value field. The inner normalization introduces no new
quantified coordinates and no function of a newly constructed linear
combination.

\subsection{Comparison of compound values}

For \(P=(\lambda,\mu)\), put \(\nu=\mu+\omega\). On a tail
satisfying the bound \(x\ge N(\nu)\), define \(K_P(x)=x\) if \(\lambda=\nu\),
and \(K_P(x)=E_{\lambda,\nu}(x)\) otherwise. The last interval
has exactly one event, so the cocycle gives
\begin{equation}\label{eq:final-peel}
 E_P(x)=f(K_P(x)).
\end{equation}

\begin{lemma}[Comparison preserving input order]\label{lem:field-comparison}
On each fixed order of the raw coordinates, every comparison of ordinary
value fields has an effective finite Boolean form in which the greatest
raw coordinate occurs only linearly. Nonlinear parameters in that form
are fields on strictly smaller coordinates.
\end{lemma}
\begin{proof}
Cancel common outer applications of \(f\) by injectivity. If an explicit
iterate remains on one side and a primitive entry on the other, peel one
last \(f\) from the latter by \eqref{eq:final-peel} and cancel again.
Each such step decreases the remaining fixed iterate count. If a primitive
becomes the identity convention, retain the resulting comparison with a
raw coordinate. Thus finitely many steps leave two primitive entries, or
a field compared with a raw coordinate. Only finitely many endpoint
advances \(\mu+\omega j\) occur.

Identify equal raw coordinates in the chosen order branch. For primitive
entries at \(x<y\), Lemma~\ref{lem:gap-input} fixes the answer if the
larger-input map has at least the larger gap. Otherwise the inverse
threshold table makes it a comparison of \(y\) with an entry at \(x\),
with its exact equality condition. At equal inputs use the gap and
equal-gap certificates. A field at the larger input is greater than that
input and hence greater than a smaller bare coordinate. A larger bare
coordinate compared with a lower-input field is already linear in the
larger coordinate. Identity versus identity is raw order. These cases
are exhaustive, and all new source and certificate domains are retained.
\end{proof}

\subsection{Annotating a count by its full leading value}

A short interval \(\lambda=\mu+\omega k\) has constant count \(k\).
For every other proper interval, put
\(\eta=\mu+\omega^2\). Let \(P'=(\lambda,\eta)\), using the
identity convention if \(\lambda=\eta\). Its actual count satisfies
\begin{equation}\label{eq:count-core}
 \tau_P(x)=\tfrac12 C_P(x)+\tau_{P'}(x)+\kappa_\mu,
 \quad C_P(x)=f(E_{P'}(x)),\quad
 \kappa_\mu=1-\lceil N(\mu)/2\rceil.
\end{equation}
This is the exact \(\omega^2\)-funnel formula, including the
\(\omega\)-tail of \(\mu\) absorbed by \(\eta\).
For \(\mu=0\) it reads \(D_{\omega^2}(q)=f(q)/2+1\), after the
preceding entry prefix. In every case, if \(q=E_{P'}(x)\) and
\(r=\tau_{P'}(x)\), then on the genuine tail
\begin{equation}\label{eq:core-error}
 q\ge x,\qquad C_P(x)\ge f(x),\qquad
 0\le r\le q\le\log_2 C_P(x).
\end{equation}

We attach \((C_P,P',\kappa_\mu)\) as an annotation to the whole count
occurrence. We do not expand it into a new independent half-value and
an active prefix count. In a class sharing the full value \(C\), ordinary
coefficient \(b\) and count coefficients \(a_i\) have combined leading
coefficient \(b+\frac12\sum_i a_i\). Its zero test is the integer test
\(2b+\sum_i a_i=0\). Treating \(C\) and \(C/2\) as separated
scales would give an incorrect test.

\begin{lemma}[Exact cancellation and transfer]\label{lem:count-transfer}
A class with zero combined leading coefficient can be replaced exactly by
its prefix counts and constants. Equal-input prefixes merge. Prefixes on
distinct inputs can be transferred to smaller input ranks whenever their
cores coincide.
\end{lemma}
\begin{proof}
For a zero class, \eqref{eq:count-core} gives the exact replacement
\begin{equation}\label{eq:zero-replacement}
 bC+\sum_i a_i\tau_{P_i}(x_i)
   =\sum_i a_i\tau_{P_i'}(x_i)+\sum_i a_i\kappa_i.
\end{equation}
All ordinary value occurrences in the class disappear. Identity and
short prefixes become constants. If two remaining prefixes have the
same input \(x\), equality of their cores gives equality of their entry
norms. These are \(f^r(x)\) and \(f^s(x)\); strict increase above the
identity gives \(r=s\). Their counts therefore merge on this guard.

Now take two remaining prefixes at \(x<y\), with
\(E_{P'}(y)=E_{Q'}(x)\). They are proper after the preceding removals.
Input comparison forbids \(\delta_{P'}\ge\delta_{Q'}\).
The complete collision table therefore forces its equality certificate
and
\[
 y=E_H(x),\qquad
 H=(\lambda_{Q'},\mu_{Q'}+\delta_{P'}).
\]
The actual \(Q'\)-cocycle factors through \(H\). At \(y\) its
remaining same-gap entry has the same terminal norm as \(P'\), hence
the same count by injectivity of iteration. Consequently
\begin{equation}\label{eq:rank-transfer}
 \tau_{P'}(y)=\tau_{Q'}(x)-\tau_H(x).
\end{equation}
This is an integer identity on the actual collision guard, even with
signed coefficients. A nonequality certificate makes the collision
branch impossible. Both terms on the right have smaller raw input rank.
Choose the smallest surviving input and transfer every higher one to it.
All intermediate target norm bounds and comparison domains are imposed.
\end{proof}

\subsection{A finite charge and the order normalization loop}

Fix a strict order of the \(r\) distinct raw coordinates, giving ranks
\(0,\ldots,r-1\). For one row put
\begin{equation}\label{eq:finite-charge}
 \mathcal M=\sum_{\text{active count occurrences}}3^{\operatorname{rank}(x)},
 \qquad
 \mathcal V=\#\{\text{ordinary nonlinear value occurrences}\}.
\end{equation}
Combine identical syntactic fields and remove zero coefficients.
Annotations add no ordinary occurrences. Pattern guards are kept
separately; Lemma~\ref{lem:field-comparison} already makes them linear
in the greatest coordinate.

\begin{lemma}[Finite inner normalization]\label{lem:inner-termination}
Every exact zero-class replacement strictly decreases
\(\mathcal M+\mathcal V\), and neither summand increases.
\end{lemma}
\begin{proof}
An ended prefix removes a count occurrence. Merging two prefixes at
one input replaces at least two equal-rank occurrences by at most one.
A transferred occurrence of rank \(j\) creates at most two lower-rank
occurrences, of total charge at most
\(2\cdot3^{j-1}<3^j\). If there is just one surviving prefix and
its count charge has not decreased, the equation \(2b+a=0\) with
\(a\ne0\) forces \(b\ne0\); an ordinary value occurrence has
been consumed. A class consisting only of ordinary values also decreases
\(\mathcal V\). Every zero class with two or more count occurrences
has a repeated input, distinct inputs, or an ended prefix, so one of
the preceding cases applies. Formula \eqref{eq:zero-replacement} creates
no ordinary values. This proves the assertion.
\end{proof}

Let \(z\) be the greatest raw coordinate. Normalize an order atom by
the following finite procedure. Remove short counts and zero coefficients.
If \(z\) has no remaining nonlinear occurrence, retain the entire row
as a linear row in \(z\), with its correlated lower-input fields as
parameters, and stop. Otherwise annotate all remaining counts, branch on
the weak order of their full cores and the ordinary nonlinear fields,
and retain every exact pattern guard. Process the largest current class.
A zero combined coefficient is canceled by Lemma~\ref{lem:count-transfer},
after which the procedure repeats. At a nonzero coefficient it stops by
the following domination argument.

Because the initial exit check failed, a nonlinear field at \(z\)
remains, so the largest current full value \(C\) satisfies \(C\ge f(z)\).
Every raw coordinate is at most \(z\le\log_2 C\). Distinct ordinary
full values are separated by adjacent growth; every smaller one is at
most \(\log_2 C\). Count errors satisfy \eqref{eq:core-error}.
The entire row is consequently
\[
 \beta C+E,\qquad 0\ne\beta\in\tfrac12\mathbb Z,
 \qquad |E|\le K\log_2 C+K_0,
\]
for effective fixed integers \(K,K_0\). An elementary cutoff makes
\(|E|<C/2\), so its sign is the sign of \(\beta\), and equality is
false. Previously canceled errors have not been dropped: they were
replaced exactly by active prefixes before this new largest class was
chosen. If those cancellations removed all nonlinear occurrences of
\(z\), the exit check instead retained its exact linear row.

The charge lemma bounds the depth of every inner branch by the initial
\(\mathcal M+\mathcal V\). Each node has finite branching. Thus the
entire symbolic normalization tree can be generated before numerical
tails are chosen. It contains finitely many interval, certificate,
peeling and transfer labels, and finitely many coefficient bounds.

\subsection{Residues, finite strips, and one-variable projection}

Proposition~\ref{prop:entry-residues} converts every current field
congruence modulo \(m\) into a finite Presburger condition on its
original raw input modulo \(2m\). For example a long count has residue
\(p(x)-c_\mu\), while a short entry has residue \(2p(x)-2k\).
These are tables in \(x\bmod2m\). All occurrences of the same \(x\)
use one shared residue. An optional explicit half-field can either use
the corresponding table or its exact natural graph \(2h=f(q)\).
The count annotations themselves introduce no such half-field atom.

For a matrix in \(n\) coordinates, first generate all finite inner
normalization trees and all current modular tables. Choose a common
cutoff \(B\) above their source norms, affordability and comparison
thresholds, prefix cutoffs, and domination margins. This maximum is
effective because the symbolic closure is finite. It does not require
finding an unknown first divergence or computing a descent length.

If some coordinate is below \(B\), branch over each of its finitely
many permitted values and substitute that same value throughout the
matrix. All fields on it are now closed exact calculations, with their
true definedness flags. The other coordinates stay unbounded: these
branches cover whole infinite strips, not merely a finite cube. Each
has one fewer coordinate and retains every earlier domain guard.

On the all-large branch, identify equal coordinates or choose a strict
base order. The order normalization and every pattern guard are linear
in the greatest coordinate \(z\). Current congruences are ordinary
Presburger congruences in the raw coordinates. Apply effective Presburger
elimination to \(\exists z\), treating all lower-input fields as
parameters \cite{cooper}. Its equivalence holds for every parameter
assignment; substituting the actual correlated values preserves it.
The resulting formula is in the same matrix language on \(n-1\)
coordinates: coefficients and moduli may change, but arguments of fields
do not change.

New moduli are handled at the next, smaller-dimensional stage, with its
own finite tails and complete small-input strips. There is no loop
raising cutoffs at the same dimension. At zero coordinates use the
exact closed evaluator. Induction on \(n\), together with the inner
charge, proves termination and exactness. A universal block is the
absence of an existential counterexample on the same original domain.
Unique circuit extensions justify the initial flattening. This proves
Theorem~\ref{thm:arithmetic}.

\subsection{Two checks on the arithmetic mechanism}

The identities
\[
 D_{\omega^2}(t)=f(t)/2+1,\qquad
 2D_{\omega^2}(t)-f(t)-1=1
\]
illustrate why full values, rather than literal half-values, must be
grouped. The leading coefficient cancels, leaving the exact constant.
For a transfer between distinct input ranks, put
\(\mu=\omega^2\), \(\lambda=\omega^2\cdot2\), and \(y=f(x)\).
The first move from \(\lambda+\omega\) reaches the \(\lambda\)-cut,
so
\[
 D_{\lambda+\omega}(x)=1+D_\lambda(y).
\]
The two long-count prefixes have equal cores. Their collision transfers
the higher-input count to the lower input with an exact difference of
one. These are symbolic identities for the genuine control; no smaller
numerical surrogate is used.

The theorem concerns closed blocks. It does not establish quantifier
elimination with retained free parameters in the same forward language.
Indeed the unary quantifier-free formulas of the retained matrix language
are eventually periodic, by the same order normalization and residue tables. The projection
\(\exists z\,[x=f(z)]\) is an infinite set with unbounded gaps and
is not eventually periodic. Thus a general projection theorem would
need a richer output language. This distinction does not obstruct the
closed counterexample sentences used for canonical comparison.
